\chapter{Background and Literature Review}
\label{sec:background}

This chapter establishes the theoretical foundations and contextualizes the relevant prior research that motivates the proposed task-conditioned latent predictive model. We begin by reviewing advancements in goal-conditioned reinforcement learning and language-guided control, followed by the evolution of model-based reinforcement learning and world models. We then detail the mathematical framework of Energy-Based Models (EBMs) and trace the recent paradigm shift toward Joint-Embedding Predictive Architectures (JEPAs) and non-contrastive self-supervised learning, focusing on theoretical mechanisms to prevent representational collapse. Finally, we survey the architectural primitives that facilitate high-dimensional representation learning.

\section{Goal-Conditioned and Language-Guided Robot Learning}

Traditional reinforcement learning (RL) formulates the control problem as the maximization of a singular, pre-defined scalar reward signal. While successful in narrow domains, this paradigm scales poorly to unstructured environments requiring multi-task generalization. Goal-Conditioned Reinforcement Learning (GCRL) addresses this limitation by parameterizing the value function and policy over a distribution of goals~\cite{liu2022goal}. A foundational step in this direction was the introduction of Universal Value Function Approximators (UVFAs)~\cite{uvfa}, which decoupled state and goal embeddings to achieve generalization over unseen targets. To mitigate the sparse reward problem inherent in GCRL, Hindsight Experience Replay (HER)~\cite{her-rl} enabled sample-efficient learning by retroactively relabeling failed trajectories as successful executions of alternative goals.

To deploy GCRL in real-world robotics, reliance on explicit ground-truth state information is often impractical, necessitating a shift toward visual reinforcement learning. Frameworks for unsupervised learning from raw pixels, such as the Unsupervised Reinforcement Learning Benchmark (URLB)~\cite{urlb}, evaluate agents based on their ability to learn behavioral priors without external reward signals. Reward-free exploration algorithms construct environment models or state-visitation distributions that can subsequently be fine-tuned for downstream tasks~\cite{rewfree-rl, rfrl-linapp}. Within the visual domain, approaches have evolved from predicting goal images directly~\cite{nair2018visual, nair2020goal} to leveraging imagined visual affordances~\cite{khazatsky2021can} and employing object-agnostic masking for generic visual goal formulation~\cite{shahriar2025general}.

Concurrently, natural language has emerged as a powerful, composable abstraction for policy conditioning and hierarchical control~\cite{luketina2019survey, jiang2019language}. Language enables agents to ingest multi-task specifications without requiring exact goal-image prototypes. Research incorporating textual advice to augment multi-goal RL (e.g., ACTRCE~\cite{actrce}) has demonstrated that language can effectively shape the exploration manifold. These developments, alongside advancements in action-free pre-training from offline videos~\cite{seo2022reinforcement} and low-cost bimanual teleoperation architectures (e.g., ACT~\cite{act}), indicate a convergence toward models that map raw sensory streams to language-grounded actions.

\section{World Models and Model-Based Control}

Model-Based Reinforcement Learning (MBRL) differentiates itself from model-free methods by explicitly learning a transition dynamics model of the environment. Rather than directly mapping states to actions via a policy network, MBRL agents leverage a ``world model'' to simulate future states and evaluate candidate action sequences via planning algorithms.

Early predictive architectures, such as ALVINN~\cite{pomerleau1988alvinn}, laid the groundwork for end-to-end sensorimotor mapping, but modern world models operate almost exclusively in compressed latent spaces. Predicting raw, high-dimensional pixel arrays forces the network to allocate substantial representational capacity to high-frequency, task-irrelevant noise (e.g., background textures, stochastic lighting variations). Latent world models circumvent this by encoding visual states into a compact representation $\mathbf{e}_t$ and predicting future states through a learned transition function $\mathbf{e}_{t+1} = h_\theta(\mathbf{e}_t, a_t)$.

Prominent architectures in this space include Temporal Difference Learning for Model Predictive Control (TD-MPC)~\cite{Hansen2022tdmpc} and its successor, TD-MPC2~\cite{hansen2024tdmpc2}. These frameworks learn latent representations by jointly optimizing for reward prediction, temporal difference value targets, and forward dynamics consistency. By executing Model Predictive Control (MPC) directly within the optimized latent manifold, these models achieve state-of-the-art sample efficiency and robustness in continuous control tasks. The predictive rollout mechanism formulated in our methodology is conceptually aligned with the forward transition models of these MBRL frameworks, albeit re-framed through an energy-based, task-conditioned objective rather than explicit reward maximization.

\section{Energy-Based Models (EBMs)}

Energy-Based Models (EBMs) provide a unified theoretical framework for learning complex dependencies by assigning a scalar ``energy'' value to configurations of variables. As formalized by LeCun et al.~\cite{lecun2006tutorial}, an EBM characterizes the compatibility between an observation $x$ and a target $y$ via an energy functional $E(x,y)$. Proper inference consists of finding the $y$ that minimizes this energy: $y^\ast = \arg \min_y E(x,y)$. The energy function models an unnormalized probability distribution over the state space via the Boltzmann distribution.

Historically, optimizing EBMs presented significant computational challenges, primarily due to the intractability of calculating the partition function (the normalizing integral over all possible states). Classical training methodologies relied on approximations like Contrastive Divergence~\cite{hinton2002training}, which required costly Markov Chain Monte Carlo (MCMC) sampling to approximate the gradient of the log-partition function. However, recent developments in score-matching and diffusion-based approximations have provided more scalable mechanisms for EBM training~\cite{train-ebm}.

In the domain of robotic control, EBMs offer a flexible alternative to standard feed-forward function approximation. Soft Actor-Critic frameworks have effectively utilized deep energy-based policies to model complex, multi-modal action distributions~\cite{haarnoja2017reinforcement}. More recently, Implicit Behavioral Cloning~\cite{florence2022implicit} demonstrated that treating imitation learning as an energy minimization problem---where the EBM assigns low energy to expert state-action pairs and high energy to all other actions---significantly outperforms explicit policy regression in tasks characterized by discontinuities and multi-modality. In this thesis, the EBM paradigm is adapted to assess the temporal and semantic compatibility of sequential state transitions under language constraints.

\section{Self-Supervised Learning and the JEPA Paradigm}

Self-Supervised Learning (SSL) aims to extract rich, actionable representations from unlabeled data. Early visual SSL relied heavily on generative, reconstruction-based proxies (e.g., autoencoders or generative adversarial networks). However, pixel-level reconstruction strictly enforces the preservation of stochastic, uncontrollable details, rendering the resulting representations sub-optimal for high-level reasoning and planning.

To address this, LeCun~\cite{lecun2022path} proposed the Path Towards Autonomous Machine Intelligence (AMI), centered around the Joint-Embedding Predictive Architecture (JEPA). JEPAs abandon raw signal reconstruction. Instead, they encode an observation $x$ and a target $y$ into a shared latent space and train a predictor network to estimate the latent representation of $y$ from the latent representation of $x$ (often conditioned on a latent action or variable $z$).

This paradigm has proven highly successful across modalities. In the static domain, the Image-JEPA (I-JEPA)~\cite{ijepa} demonstrated that predicting the latent representations of masked image patches yields highly semantic features without relying on hand-crafted data augmentations. This was subsequently extended to the spatio-temporal domain with Video-JEPA (V-JEPA)~\cite{vjepa2}, which masks out entire continuous regions of video tubes, forcing the model to learn intuitive physics and temporal continuity. For control architectures, action-conditioned JEPAs such as ACT-JEPA~\cite{actjepa}, TD-JEPA~\cite{bagatella2025td}, and stable end-to-end frameworks like LeWorldModel~\cite{Maes2026LeWorldModelSE} have demonstrated that predicting future latents as a function of current states and actions yields representations that are highly amenable to zero-shot or sample-efficient reinforcement learning.

\subsection{Latent Collapse and Non-Contrastive Regularization}

A fundamental vulnerability of any joint-embedding architecture is the risk of \textit{latent collapse}. If the encoder maps all inputs to a constant vector, the predictor can trivially minimize the prediction error to zero, resulting in a functionally useless representation. Early SSL methods prevented collapse via contrastive learning, wherein the objective simultaneously pushed embeddings of different inputs (negative pairs) apart while pulling augmented views of the same input (positive pairs) together. However, requiring large batches of carefully curated negative samples incurs severe memory costs and scaling limitations.

Non-contrastive regularization methods eliminate the need for negative sampling by mathematically bounding the geometry of the latent manifold. An early success in this domain was VICReg (Variance-Invariance-Covariance Regularization)~\cite{bardes2021vicreg}, which explicitly forced the variance of each latent dimension to remain above a specified threshold while decorrelating the features via a covariance penalty.

More recently, Sketched Isotropic Gaussian Regularization (SIGReg), introduced as the cornerstone of the LeJEPA architecture~\cite{lejepa}, provided a provable, information-theoretic solution to latent collapse without relying on heuristic hyperparameter tuning. SIGReg enforces that the distribution of latent encodings approximates an isotropic Gaussian $\mathcal{N}(0, I)$. It achieves this efficiently via the method of ``sketching'': random one-dimensional directions $\omega$ are sampled from a unit sphere $\mathcal{S}^d$, and the high-dimensional latents are projected onto these directions. An Epps-Pulley normality test is computed on these 1D projections, heavily penalizing deviations from a standard normal distribution. By integrating SIGReg, the methodology detailed in this thesis ensures that the latent manifold retains maximal information capacity, enabling robust multi-step task-conditioned rollouts.

\section{Architectural Foundations for Latent Representation}

The realization of latent world models and EBMs is deeply tied to the underlying neural architectures. For visual encoding, the field has transitioned from Convolutional Neural Networks (CNNs), such as the ResNet family~\cite{DBLP:journals/corr/HeZRS15}, toward Vision Transformers (ViTs)~\cite{vit}. ViTs process images as sequences of flattened patches, utilizing multi-head self-attention to capture global receptive fields natively. Foundation models leveraging the ViT architecture, such as DINOv3~\cite{dinov3}, produce generic, dense visual features that exhibit strong part-level understanding and physical tracking capabilities without explicit semantic supervision.

For sequence modeling and predictive mapping, the standard Transformer architecture~\cite{vaswani2017attention} remains dominant. Modern instantiations frequently replace standard feed-forward layers with SwiGLU (Swish-Gated Linear Units)~\cite{swiglu}, which have been shown to improve gradient flow and representation capacity in deep generative models.

Finally, to bridge the modality gap between language instructions and visual observations, multimodal foundation models like CLIP (Contrastive Language-Image Pre-training)~\cite{clip} are utilized. CLIP maps text strings and images into a shared semantic latent space, making it an ideal off-the-shelf encoder for goal-conditioning. However, because foundational text encoders are generally trained on static web data, their representations often lack the fine-grained physical nuance required for manipulation tasks. To adapt these models without incurring catastrophic forgetting or requiring the computational overhead of full fine-tuning, Parameter-Efficient Fine-Tuning (PEFT) methods are employed. Specifically, Low-Rank Adaptation (LoRA)~\cite{hu2022lora} injects trainable rank-decomposition matrices into the Transformer layers, allowing the linguistic encoder to align with the robotic visual manifold while keeping the majority of the pre-trained weights frozen.

\nomenclature{RL}{Reinforcement Learning}
\nomenclature{GCRL}{Goal-Conditioned Reinforcement Learning}
\nomenclature{UVFA}{Universal Value Function Approximator}
\nomenclature{HER}{Hindsight Experience Replay}
\nomenclature{URLB}{Unsupervised Reinforcement Learning Benchmark}
\nomenclature{EBM}{Energy-Based Model}
\nomenclature{MCMC}{Markov Chain Monte Carlo}
\nomenclature{SSL}{Self-Supervised Learning}
\nomenclature{AMI}{Autonomous Machine Intelligence}
\nomenclature{JEPA}{Joint-Embedding Predictive Architecture}
\nomenclature{I-JEPA}{Image Joint-Embedding Predictive Architecture}
\nomenclature{V-JEPA}{Video Joint-Embedding Predictive Architecture}
\nomenclature{VICReg}{Variance-Invariance-Covariance Regularization}
\nomenclature{SwiGLU}{Swish-Gated Linear Unit}
\nomenclature{CLIP}{Contrastive Language-Image Pre-training}
\nomenclature{LoRA}{Low-Rank Adaptation}
